\section{Feasible Branch Shifts}
\label{sec:shifts}

This section formalizes the concept of the feasible branch shift space, which characterizes the coordinate transformations that preserve active distance constraints. 
Geometrically, moving from a branch code $s$ to $s \oplus h$ via a branch shift $h \in \mathbb{F}_2^{B_{\mathrm{c}}}$ corresponds to a sequence of partial reflections. 
A shift preserves the validity of all active constraints if it is produced by a generator combination with zero net violations, meaning $\alpha \in \ker V$. 
Because multiple generator combinations can yield the same branch mask, the space of valid branch shifts $\mathcal{K}_F$ is defined as the image of this kernel under the mask matrix $M$:
\[
\mathcal{K}_F = M(\ker V).
\]
To ensure that only shifts within $\mathcal{K}_F$ preserve these distances under generic geometric configurations, we introduce the \emph{generic mirror separation} hypothesis.

\begin{definition}[Generic mirror separation]
\label{def:mirror-sep}
Let $P$ be the real algebraic set of nondegenerate lateration parameters (predecessor coordinates and active edge lengths). 
The instance satisfies generic mirror separation if there exists a proper algebraic exceptional set $Z \subset P$ such that, for all parameters in $P \setminus Z$, any shift $h \notin \mathcal{K}_F$ and branch code $s$ satisfy
\[
\|x_a(s \oplus h) - x_b(s \oplus h)\|_2 \neq \|x_a(s) - x_b(s)\|_2
\]
for at least one active edge $\{a,b\} \in F$ (where both embeddings are defined).
\end{definition}

We use generic mirror separation as an explicit instance-level hypothesis. The main theorem below does not assert that every ordered instance satisfies it; rather, it proves the rank count for instances that do.

The next section uses this hypothesis to prove the algebraic rank count.
